\subsection{弗雷德霍姆映射的摄动}

定理 1. --- 设 E 与 F 是巴拿赫空间。从 E 到 F 的弗雷德霍姆算子组成的集 \(\mathcal{F}(\mathrm{E};\mathrm{F})\) 在巴拿赫空间 \(\mathcal{L}(\mathrm{E};\mathrm{F})\) 中是开的，且映射 \(u\mapsto\mathrm{ind}(u)\)（它从 \(\mathcal{F}(\mathrm{E};\mathrm{F})\) 到 \(\mathbf{Z}\) 中）是局部常值的。

该定理由下面更精确的命题得出：

命题 4. --- 设 E 与 F 是巴拿赫空间，\(u_{0}\colon\mathrm{E}\to\mathrm{F}\) 是弗雷德霍姆算子，\(v_{0}\) 是 \(u_{0}\) 的一个拟逆。存在 \(u_{0}\) 在 \(\mathcal{L}(\mathrm{E};\mathrm{F})\) 中的一个开邻域 U 以及一个解析映射 \(\varphi\colon\mathrm{U}\to\mathcal{L}(\mathrm{F};\mathrm{E})\)，使得

\begin{enumerate}
\def\labelenumi{(\roman{enumi})}
\item
  我们有 \(\varphi(u_{0}) = v_{0}\)；
\item
  对每个 \(u\)（在 \(U\) 中），映射 \(\varphi(u)\) 是 \(u\) 的拟逆，特别地，\(u\) 是弗雷德霍姆映射；
\item
  对 U 中的每个 \(u\)，我们有 \(\operatorname{ind}(u) = \operatorname{ind}(u_0)\)。
\end{enumerate}

由 III, p. 42 的 prop. 2 的 (ii)，存在拓扑直和分解 \(\mathrm{E} = \mathrm{E}_1 \oplus \mathrm{E}_2\) 与 \(\mathrm{F} = \mathrm{F}_1 \oplus \mathrm{F}_2\)，且存在 \(\alpha_0 \in \mathcal{I}(\mathrm{E}_1; \mathrm{F}_1)\)，使得 \(\mathrm{E}_2\) 与 \(\mathrm{F}_2\) 是有限维的，而 \(u_0\) 相对于这些分解由矩阵 \(\left( \begin{array}{cc}\alpha_0 & 0\\ 0 & 0 \end{array} \right)\) 表示。

设 U 是由这样的元素 \(u\)（它属于 \(\mathcal{L}(\mathrm{E};\mathrm{F})\)）组成的集：在矩阵表示 \(\left( \begin{array}{cc}\alpha & \beta \\ \gamma & \delta \end{array} \right)\)（即 \(u\) 相对于这些分解的矩阵表示）中，我们有 \(\alpha \in \mathcal{I}(\mathrm{E}_1;\mathrm{F}_1)\)。由于 \(\mathcal{I}(\mathrm{E}_1;\mathrm{F}_1)\) 在 \(\mathcal{L}(\mathrm{E}_1;\mathrm{F}_1)\) 中是开的

（III, p. 57 的 prop. 3），U 是 \(u_0\) 在 \(\mathcal{L}(\mathrm{E};\mathrm{F})\) 中的开邻域。对 U 中的 \(u = \begin{pmatrix} \alpha & \beta \\ \gamma & \delta \end{pmatrix}\)，置

\[
\varphi (u) = v _ {0} + \left( \begin{array}{c c} \alpha^ {- 1} - \alpha_ {0} ^ {- 1} & 0 \\ 0 & 0 \end{array} \right).\tag{1}
\]

我们有 \(\varphi(u_{0}) = v_{0}\)，且映射 \(\varphi\) 是解析的（loc. cit.）。模有限秩连续线性映射，我们有同余关系

\[
u \equiv \left( \begin{array}{c c} \alpha & 0 \\ 0 & 0 \end{array} \right) , v _ {0} \equiv \left( \begin{array}{c c} \alpha_ {0} ^ {- 1} & 0 \\ 0 & 0 \end{array} \right) , \varphi (u) \equiv \left( \begin{array}{c c} \alpha^ {- 1} & 0 \\ 0 & 0 \end{array} \right).
\]

因此，\(\varphi(u)\) 是 \(u\) 的拟逆。U 的每个元素 \(u\) 通过限制定义了一个从 \(E _1\) 到 \(F _2\) 在 F 中的一个拓扑补上的同构。由 III, p. 44 的 prop. 3，我们因此有

\[
\begin{array}{c} \operatorname{ind} (u) = \operatorname{codim} _ {\mathrm{F}} (u (\mathrm{E} _ {1})) - \operatorname{codim} _ {\mathrm{E}} (\mathrm{E} _ {1}) \\ = \dim (\mathrm{F} _ {2}) - \dim (\mathrm{E} _ {2}) = \operatorname{ind} (u _ {0}), \end{array}
\]

证毕。
% semantic-fragments: legacy-input-seam
\relax{}
